\begin{figure}[H]
\centering
\begin{tikzpicture}[>=Stealth,font=\small]
\node at (0,1.65)
 {$(\FeatureCoordinates,\NormSymbol)$};
\node at (5.8,1.65)
 {$(\OrthogonalChange\circ\FeatureCoordinates,
     \NormSymbol\circ\OrthogonalChange^{-1})$};
\begin{scope}
 \filldraw[fill=black!4] (0,0) ellipse (1.45 and 0.88);
 \draw[->,black!35] (-1.8,0) -- (1.8,0);
 \draw[->,black!35] (0,-1.15) -- (0,1.2);
 \foreach \x/\y in {-0.8/-0.3,-0.55/0.42,0.12/-0.55,0.48/0.3,0.92/-0.15}
  \fill (\x,\y) circle (1.4pt);
 \draw[black!55] (0.48,0.3) -- (0.85,1.03);
 \node at (0.85,1.25)
  {$\FeatureCoordinates(\ComparisonPoint)$};
 \draw[black!55] (-0.8,-0.3) -- (-0.95,-0.9);
 \node at (-0.95,-1.12)
  {$\FeatureCoordinates(\IntegrationPoint)$};
\end{scope}
\begin{scope}[xshift=5.8cm]
 \begin{scope}[rotate=30]
  \filldraw[fill=black!4] (0,0) ellipse (1.45 and 0.88);
 \end{scope}
 \draw[->,black!35] (-1.8,0) -- (1.8,0);
 \draw[->,black!35] (0,-1.15) -- (0,1.2);
 \begin{scope}[rotate=30]
  \foreach \x/\y in {-0.8/-0.3,-0.55/0.42,0.12/-0.55,0.48/0.3,0.92/-0.15}
   \fill (\x,\y) circle (1.4pt);
  \coordinate (rotatedy) at (0.48,0.3);
  \coordinate (rotatedz) at (-0.8,-0.3);
 \end{scope}
 \draw[black!55] (rotatedy) -- (0.4,1.03);
 \node at (0.4,1.25)
  {$\OrthogonalChange\FeatureCoordinates(\ComparisonPoint)$};
 \draw[black!55] (rotatedz) -- (-1,-0.92);
 \node at (-1,-1.14)
  {$\OrthogonalChange\FeatureCoordinates(\IntegrationPoint)$};
\end{scope}
\draw[->] (2.05,0.15) -- node[above]
 {$\OrthogonalChange\in\OrthogonalGroup(2)$} (3.75,0.15);
\node at (0,-1.55)
 {$\{\FirstVector\mid\NormSymbol(\FirstVector)\leq1\}$};
\node at (5.8,-1.55)
 {$\{\FirstVector\mid
     (\NormSymbol\circ\OrthogonalChange^{-1})(\FirstVector)\leq1\}$};
\node at (2.9,-2.35)
 {$(\NormSymbol\circ\OrthogonalChange^{-1})
     \bigl(\OrthogonalChange\FeatureCoordinates(\ComparisonPoint)
           -\OrthogonalChange\FeatureCoordinates(\IntegrationPoint)\bigr)
   =\NormSymbol\bigl(\FeatureCoordinates(\ComparisonPoint)
                     -\FeatureCoordinates(\IntegrationPoint)\bigr)$};
\end{tikzpicture}
\caption{An orthogonal change of coordinates acts on both the coordinate
image and the norm unit ball.
The dots represent a finite coordinate image in dimension
$2$,
and the shaded regions represent the two unit balls.
Their simultaneous rotation leaves the induced pseudometric unchanged.}
\label{fig:orthogonal-coordinate-change}
\end{figure}
